\documentclass[a4paper, 14pt]{extarticle}

\usepackage{../latexDependencies/misc/preamble2}

\geometry{a4paper}

% Название дисциплины
\newcommand{\subject}{Численные методы} 

% Тип работы
% lab - для лабораторной работы 
% hw  - для домашней     работы
\newcommand{\task}{hw} 

% Номер работы
\newcommand{\taskNumber}{2-1} 

% Название работы
\newcommand{\taskNameOne}{Интерполяция Лагранжа.} 
\newcommand{\taskNameTwo}{Вычисление интерполяционного полинома} 
\newcommand{\taskNameThree}{Лагранжа.} 

% Имя студента
\newcommand{\studentName}{Очкин Н.В.}

% Имя преподававателя
\newcommand{\teacherName}{Кутыркин В.А.}

% Группа
\newcommand{\group}{ФН11-52Б}

% Вариант
\newcommand{\variant}{9}

\begin{document}

\graphicspath{ {../latexDependencies/images} } 
\normalsize

\newcommand{\printTask}{%
    \ifthenelse{\equal{\task}{lab}}{%
        лабораторной%
    }{%
        \ifthenelse{\equal{\task}{hw}}{%
            домашней%
        }{%
            Неизвестный тип задания%
        }%
    }%
}

\begin{titlepage}

    \begin{center}
        {\footnotesize \itshape Федеральное государственное бюджетное 
                       образовательное учреждение высшего образования}
    \end{center}

    \begin{minipage}[c]{0.1\textwidth}
        \includegraphics[width=1.1\textwidth]{iconBMSTU}
    \end{minipage}
    \hfill
    \begin{minipage}[c]{0.9\textwidth}
        \centering
        \itshape
        \bfseries
        \small
        \guillemotleft Московский государственный технический университет \\
        имени Н.Э. Баумана\guillemotright \\
        (национальный исследовательский университет) \\
        (МГТУ им. Н.Э. Баумана) 
    \end{minipage}

    \vspace{0.5cm}
    \noindent\rule{\textwidth}{2pt} \\

    \noindent\uline{\textbf{ФАКУЛЬТЕТ} ФУНДАМЕНТАЛЬНЫЕ НАУКИ} \\
    \vspace{-5pt} \\
    \noindent\uline{\textbf{КАФЕДРА} ВЫЧИСЛИТЕЛЬНАЯ МАТЕМАТИКА И МАТЕМАТИЧЕСКАЯ} \\
    \vspace{-5pt} \\
    \noindent\uline{ФИЗИКА (ФН11)} \\
    \vspace{-5pt} \\
    \noindent\uline{\textbf{НАПРАВЛЕНИЕ ПОДГОТОВКИ} МАТЕМАТИКА И КОМПЬЮТЕРНЫЕ} \\
    \vspace{-5pt} \\
    \noindent\uline{НАУКИ (02.03.01)} \\

    \begin{center}
        \bfseries
        \textsc{О т ч е т} \\[10pt]
        по \printTask {} работе \textnumero {} \taskNumber
    \end{center}

    \vspace{10pt}

    \hspace{10pt} 
    \noindent \textbf{Название \printTask {} работы:} \par
    \vspace{5pt}
    \hspace{10pt} 
    \noindent \textbf{\uline{\taskNameOne}} \vspace{5pt} \\
    \null\hspace{31pt} 
    \textbf{\uline{\taskNameTwo}} \vspace{5pt} \\
    \null\hspace{31pt} 
    \textbf{\uline{\taskNameThree}}

    \vspace{10pt}

    \begin{center}
        \bfseries
        Вариант \textnumero {} \variant
    \end{center}

    \vspace{20pt}

    \hspace{10pt} 
    \noindent \textbf{Дисциплина:} \par
    \vspace{10pt}
    \hspace{10pt} 
    \noindent {\large \subject}

    \vspace{10pt}

    \begin{flushright}
        \renewcommand{\arraystretch}{3}
        \begin{tabular}{r r r}
            \multicolumn{1}{l}{Студент группы \uline{\group}} & 
            $\quad \underset{\text{(Подпись, дата)}}{\underline{\hspace{3cm}}} \quad$ & 
            \multicolumn{1}{c}{$\underset{\text{(И.О. Фамилия)}}{\uline{\textbf{\studentName}}}$} \\

            \multicolumn{1}{l}{Преподаватель} & 
            $\quad \underset{\text{(Подпись, дата)}}{\underline{\hspace{3cm}}} \quad$ & 
            \multicolumn{1}{c}{$\underset{\text{(И.О. Фамилия)}}{\uline{\textbf{\teacherName}}}$} \\
        \end{tabular}
    \end{flushright}

    \vfill

    \begin{center}
        \small
        Москва, 2024
    \end{center}
\end{titlepage}


\newgeometry{left=25mm, right=25mm, top=20mm, bottom=20mm}

\graphicspath{ {../latexDependencies/images/HW2-1} }

% Customize section, subsection, subsubsection and paragraph styles
\titleformat{\section}
  {\normalfont\large\bfseries}{\thesection}{1em}{}

\titleformat{\subsection}
  {\normalfont\normalsize\bfseries}{\thesubsection}{1em}{}

\titleformat{\subsubsection}
  {\normalfont\small\bfseries}{\thesubsubsection}{1em}{}

\titleformat{\paragraph}
  {\small\small\bfseries}{\theparagraph}{1em}{}

% \thispagestyle{empty}

% \null\newpage

% \setcounter{tocdepth}{5}
% \setcounter{secnumdepth}{5}

% \pagenumbering{roman}

% \tableofcontents
% \newpage

% \pagenumbering{arabic}
% \setcounter{page}{1}

\setstretch{1}
\linespread{1.1}

\setlength{\parindent}{0pt}

\fontsize{12pt}{16pt}\selectfont

% \definecolor{myblue}{HTML}{0A88C2}
% \definecolor{myred}{HTML}{FF1B1C}
% \definecolor{mygreen}{HTML}{386641}

% \lstdefinestyle{mystyle}{
%     basicstyle=\ttfamily\footnotesize,
%     keywordstyle=\color{myblue},
%     stringstyle=\color{myred},
%     commentstyle=\color{green!50!black},
%     showstringspaces=false,
%     frame=leftline, 
%     framesep=10pt, 
% }

% % Set the style for Python code
% \lstset{style=mystyle, extendedchars=\true}

% --------------------------------------START--------------------------------------

\section*{Задание}\vspace{-20pt}\rule{\linewidth}{0.1mm}

\begin{equation*}
    \begin{cases}
        a_0 + b_0 (x - x_0) + c_0 (x - x_0)^2 + d_0 (x - x_0)^3, & x \in [x_0, x_1] \\ 
        a_1 + b_1 (x - x_1) + c_1 (x - x_1)^2 + d_1 (x - x_1)^3, & x \in [x_1, x_2] \\ 
        \vdots \\ 
        a_{n-2} + b_{n-2} (x - x_{n-2}) + c_{n-2} (x - x_{n-2})^2 + d_{n-2} (x - x_{n-2})^3, & x \in [x_{n-2}, x_{n-1}] 
    \end{cases}   
\end{equation*}

Таким образом, для записи функции естественного кубического сплайна дефекта 1 вам нужно просто подставить найденные коэффициенты $ a_i, b_i, c_i, d_i $ в 
вышеуказанную формулу для каждого интервала $[x_i, x_{i+1}]$.

\end{document}